% 来源：OEBPS/Text/01_34.xhtml（OPF spine 顺序；仅第一册）
\cleardoublepage
\BookEntry[菲利普·A.费舍《怎样选择成长股》]{菲利普·A.费舍《怎样选择成长股》}

\BookCover{images/书单19怎样选择成长股.jpg}{[美] 菲利普·A. 费舍著| 罗耀宗译| 地震出版社2013年版}

\begin{BookSummaryBox}
\textbf{广受华尔街推崇和敬重，以增长为导向的投资方法的先驱之作}
\par
\end{BookSummaryBox}

\subsubsection*{这本书谈什么？}

本书分三个部分阐释了以增长为导向的投资方法。第一部分是“怎样选择成长股”，包括调查上市公司的“闲聊法”、寻找优良普通股的“15个要点”、何时买进卖出、投资人“五不”原则和投资人“另五不”原则等内容。第二部分具体分析了保守型投资的四个要素，告诉新老投资者如何长期抱牢成长股而不必冒太大的风险。第三部分是“发展投资哲学”，详述了作者如何从半个世纪的商业经验中发展出自己独特投资哲学的过程。
\par

菲利普·A.费舍是广受华尔街推崇和敬重的投资家，他在40年前提出自己的投资哲学，如今依然被多数人奉为投资理财圣经。
\par

\subsubsection*{芒格为什么谈起这本书？}

芒格说及如何选股时，称：
\par

\begin{quotation}
沃伦在商学院上课时，对学生们这样讲：“我用一张考勤卡就能改善你最终的财务状况；这张卡片上有20格，所以你只能有20次打卡的机会——这代表你一生中所能拥有的投资次数。当你把卡打完之后，你就再也不能进行投资了。在这样的规则之下，你将会真正地慎重考虑你做的事情，你将不得不花大笔资金在你真正想投资的项目上。这样你的表现将会好得多。”
\par
\end{quotation}

而关于如何投资，芒格坦陈自己很认同菲利普·A.费舍在《怎样选择成长股》中的观点，即选择那些能够“增长，增长，再增长”的高品质公司的股票。费舍是关注成长股的先驱，这种思想是价值投资的一个变种和分支。在半个世纪后的今天，这种方法依然是投资界的主流投资模式之一，实际上也折射出一个极富远见的思想——少即是多。
\par

\subsubsection*{《怎样选择成长股》作者菲利普·A.费舍如是说}

“投资人‘五不’原则：
\par

1. 不买处于创业阶段的公司。
\par

2. 不要因为一只好股票在‘店头市场’交易，就弃之不顾。
\par

3. 不要因为你喜欢某公司年报的‘格调’，就去买该公司的股票。
\par

4. 不要以为一公司的本益比高，便表示未来的盈余成长已大致反映在价格上。
\par

5. 不要锱铢必较。”
\par
